\subsection{Common form of the shared-site correction}
\label{sec:shared-site-common-form}

The preceding construction defines the continuum source $\mathcal C$ without
assuming a particular lattice model.  We now give a common expression for
this source under the scaling realized by the two chains below.  This step
uses the spin-$\tfrac12$ local operator algebra and the lattice Sutherland
relation; it is not a new statement about the general nonmultiplicativity of
coherent-state symbols.

Let $P$ denote the permutation operator on two physical spin-$\tfrac12$
spaces.  After the model-dependent normalization and spectral scaling, assume
that the spatial Lax operator, the two-site Hamiltonian and the spectral
derivative have expansions
\begin{align}
\widehat L
&=\Id+\epsilon X+\epsilon^2Y+\mathcal O(\epsilon^3),
\nonumber\\
h
&=2P+\epsilon h^{(1)}+\mathcal O(\epsilon^2),
\nonumber\\
\partial_u\widehat L
&=\epsilon^2 Z+\mathcal O(\epsilon^3).
\label{eq:common-operator-expansions}
\end{align}
Here $X,Y,Z$ act on the auxiliary space and one physical site, while
$h^{(1)}$ acts on two neighbouring physical sites.  The labels $1,2$ below
specify the physical site on which a one-site operator acts; the auxiliary
space is shared and the order of auxiliary-space matrix products is retained.
For the scalings of Sections~\ref{sec:class5-model} and
\ref{sec:class6-model}, $Z$ is the coefficient obtained by differentiating
with respect to the quantum spectral parameter at fixed lattice coupling.

The order-$\epsilon^2$ coefficient of the Sutherland relation is conveniently
measured by
\begin{equation}
\mathcal B
:=[2P,X_2X_1+Y_2+Y_1]
+[h^{(1)},X_2+X_1]-Z_1+Z_2.
\label{eq:common-Sutherland-residual}
\end{equation}
For an integrable lattice Lax operator obeying the Sutherland relation at this
order, $\mathcal B=0$.  We also use the symmetric two-spin projector
$P_+:=(\Id+P)/2$ and assume
\begin{equation}
P_+h^{(1)}P_+=cP_+
\label{eq:common-triplet-condition}
\end{equation}
for a field-independent scalar $c$.  This condition holds for the two chains
studied below.  It is a sufficient condition for the common identity derived
here; no claim is made that it is necessary for arbitrary spin chains.

At a fixed continuum spin $\boldsymbol S$, define
\begin{equation}
U:=X^\downarrow,
\qquad
\Xi:=(X^2)^\downarrow-U^2,
\label{eq:common-Xi-definition}
\end{equation}
where the operator square is formed before taking the lower symbol.  Expanding
the finite-lattice time operator gives
\begin{align}
A^{(1)}&=-\ii[2P,X_2],
\nonumber\\
A^{(2)}&=\ii X_2[2P,X_2]-\ii[2P,Y_2]
-\ii[h^{(1)},X_2]-\ii Z_2.
\label{eq:common-time-coefficients}
\end{align}
The order-$\epsilon^2$ parts of the two overlapping products in
Eq.~\eqref{eq:exact-lower-symbol-lattice-zc} have equal coincident-spin
values,
\begin{equation}
\mathcal D^{(+)}_2(\boldsymbol S,\boldsymbol S)
=\mathcal D^{(-)}_2(\boldsymbol S,\boldsymbol S)
=2\ii\Xi.
\label{eq:common-quadratic-overlap}
\end{equation}
They therefore cancel before the continuum limit is taken.  Their first
spatial Taylor coefficients instead add and give
\begin{equation}
\partial_x S^i
\left.
\left(
\frac{\partial\mathcal D^{(+)}_2}{\partial S_R^i}
+\frac{\partial\mathcal D^{(-)}_2}{\partial S_L^i}
\right)\right|_{\boldsymbol S_L=\boldsymbol S_R=\boldsymbol S}
=2\ii\partial_x\Xi.
\label{eq:common-quadratic-Taylor}
\end{equation}
The coincident-spin cubic terms give the complementary contribution
\begin{equation}
\mathcal D^{(+)}_3(\boldsymbol S,\boldsymbol S)
-\mathcal D^{(-)}_3(\boldsymbol S,\boldsymbol S)
=2\ii[\Xi,U]-\ii\langle X_1\mathcal B\rangle.
\label{eq:common-cubic-overlap}
\end{equation}
Here $\langle\cdots\rangle$ denotes the lower symbol on the two physical
sites at the same spin, with auxiliary-space indices retained.  The terms
containing $Y$ are required in the intermediate expansion but cancel after
all contributions in Eq.~\eqref{eq:common-time-coefficients} are combined.

Equations~\eqref{eq:common-quadratic-Taylor} and
\eqref{eq:common-cubic-overlap} give
\begin{equation}
\mathcal C
=2\ii\bigl(\partial_x\Xi+[\Xi,U]\bigr)
-\ii\langle X_1\mathcal B\rangle.
\label{eq:common-source-with-residual}
\end{equation}
For the lattice models considered here, the Sutherland relation gives
$\mathcal B=0$, and hence
\begin{equation}
\mathcal C
=2\ii\bigl(\partial_x\Xi+[\Xi,U]\bigr).
\label{eq:common-source-general}
\end{equation}
The source is then absorbed by the local correction
\begin{equation}
\Delta V=2\ii\Xi+f(\lambda)\Id,
\label{eq:common-DeltaV-general}
\end{equation}
where $f(\lambda)$ is independent of the spin and of the spacetime
coordinates.  The model calculations in Sections~\ref{sec:class5-time-lax}
and \ref{sec:class6-time-lax} fix convenient representatives and independently
verify every step of the continuum limit.

For spin $\tfrac12$, the second moment in
Eq.~\eqref{eq:common-Xi-definition} is also determined by the spin dependence
of $U$.  Writing
\begin{equation}
X=\mathsf X_0\otimes\Id_2+\mathsf X_i\otimes\sigma^i,
\qquad
U=\mathsf X_0+\mathsf X_i S^i,
\label{eq:common-X-Pauli-expansion}
\end{equation}
one obtains
\begin{equation}
\Xi
=\mathsf X_i\mathsf X_j
\left(\delta_{ij}-S^iS^j+\ii\varepsilon_{ijk}S^k\right).
\label{eq:common-Xi-from-U}
\end{equation}
Thus a higher coefficient of the spatial-operator expansion can provide a
useful model-specific representation of $\Xi$, as happens in Class 6, without
being independent data in the common formula.
